\documentclass[aps,prl,twocolumn,superscriptaddress,nofootinbib]{revtex4-2}

\usepackage{amsmath,amssymb}
\usepackage{graphicx}
\usepackage{hyperref}
\usepackage{xcolor}
\usepackage{booktabs}
\usepackage{siunitx}

% SI units shortcuts for cQED
\DeclareSIUnit\GHz{\giga\hertz}
\DeclareSIUnit\MHz{\mega\hertz}
\DeclareSIUnit\kHz{\kilo\hertz}
\DeclareSIUnit\ns{\nano\second}
\DeclareSIUnit\us{\micro\second}

\begin{document}

\title{STUDY TITLE}
\author{Author Name}
\affiliation{Shyam Shankar Quantum Circuits Group}
\date{\today}

\begin{abstract}
% Brief summary of the study: system, method, key result.
\end{abstract}

\maketitle

% ===================================================================
\section{Motivation}
% ===================================================================
% What was studied and why.
% For REP class: cite the original paper being reproduced.
% For OPT class: describe the optimization target and its significance.


% ===================================================================
\section{System and Methods}
% ===================================================================

\subsection{Hamiltonian}
% Write out the full Hamiltonian in the appropriate frame.
% Example:
% \begin{equation}
%   \hat{H} / \hbar = \omega_q \hat{q}^\dagger \hat{q}
%                    + \frac{\alpha}{2} \hat{q}^\dagger \hat{q}^\dagger \hat{q} \hat{q}
%                    + \omega_s \hat{s}^\dagger \hat{s}
%                    + \chi \hat{q}^\dagger \hat{q} \hat{s}^\dagger \hat{s}
% \end{equation}

\subsection{Parameters}
% Table of all system parameters.
% \begin{table}[h]
% \caption{System parameters.\label{tab:params}}
% \begin{ruledtabular}
% \begin{tabular}{lcc}
%   Parameter & Symbol & Value \\
%   \hline
%   Qubit frequency    & $\omega_q/2\pi$ & \SI{5.0}{\GHz} \\
%   Anharmonicity      & $\alpha/2\pi$   & \SI{-200}{\MHz} \\
%   Storage frequency   & $\omega_s/2\pi$ & \SI{7.5}{\GHz} \\
%   Dispersive shift    & $\chi/2\pi$     & \SI{-1.5}{\MHz} \\
% \end{tabular}
% \end{ruledtabular}
% \end{table}

\subsection{Simulation Approach}
% Describe the simulation method:
% - Which cqed_sim model class was used
% - Pulse sequences and drive configurations
% - Hilbert space truncation dimensions
% - Open vs. closed system (noise model if applicable)
% - Solver settings (nsteps, tolerances)


% ===================================================================
\section{Results}
% ===================================================================
% Present quantitative findings with figures and tables.
%
% Include figures:
% \begin{figure}[t]
%   \includegraphics[width=\columnwidth]{../figures/FIGURE_NAME.pdf}
%   \caption{Description.\label{fig:label}}
% \end{figure}
%
% Reference: As shown in Fig.~\ref{fig:label}, ...


% ===================================================================
\section{Discussion}
% ===================================================================
% Physical interpretation of results.
% Comparison to literature (percent error, fidelity agreement).
% Limitations and potential sources of discrepancy.
% Convergence summary.


% ===================================================================
\section{Conclusion}
% ===================================================================
% Summary of findings.
% Potential follow-up work or open questions.


% ===================================================================
% References
% ===================================================================
\bibliographystyle{apsrev4-2}
\bibliography{references}

\end{document}
